\chapter{Dietitians}

\section*{Definition and Overview}
A dietitian is a health professional with university training in nutrition who helps people understand food and make dietary changes to improve health. Accredited practising dietitians (APDs) in Australia are qualified to provide medical nutrition therapy for conditions such as diabetes, heart disease, kidney disease, and food allergies, as well as general healthy eating advice and weight management. Dietitians translate nutrition science into practical advice and work with people of all ages. This chapter explains what dietitians do, when to see one, and how to find one.

\section*{Epidemiology}
Diet-related conditions are among the most common health problems in Australia, including obesity, type 2 diabetes, heart disease, and malnutrition. Many Australians could benefit from professional nutrition advice, but dietitian services are under-used, partly because of cost and awareness. Dietitians work in hospitals, community health, private practice, aged care, and public health. Medicare rebates are available for some dietitian services through chronic disease management plans, and private health insurance may cover others.

\section*{Aetiology and Pathophysiology}
\begin{itemize}
    \item \textbf{Training:} dietitians complete an accredited university degree and ongoing professional development
    \item \textbf{Assessment:} they assess dietary intake, nutritional status, and health conditions
    \item \textbf{Medical nutrition therapy:} tailored dietary treatment for medical conditions, such as adjusting carbohydrate for diabetes or potassium for kidney disease
    \item \textbf{Individualised advice:} practical, personalised plans that fit the person's culture, budget, and lifestyle
    \item \textbf{Education:} teaching people how to read labels, plan meals, and make sustainable changes
    \item \textbf{Scope:} dietitians differ from nutritionists, who may not have the same clinical training or accreditation
\end{itemize}

\section*{Clinical Features}
\begin{itemize}
    \item A dietitian consultation involves a detailed review of eating habits, health history, and goals
    \item Advice is practical and individualised, not a generic diet
    \item They provide support for weight management, diabetes, heart health, gut problems, and food allergies
    \item They help with tube feeding and special diets in hospital settings
    \item They can work with the whole family on healthy eating
    \item Follow-up sessions monitor progress and adjust the plan
\end{itemize}

\section*{Differential Diagnosis}
\begin{itemize}
    \item Accredited practising dietitian versus nutritionist, who may have different qualifications
    \item Dietitian advice versus general online information and fad diets
    \item Individual nutrition therapy versus population dietary guidelines
\end{itemize}

\section*{When to Seek Medical Care}
See a dietitian for help managing a medical condition through diet, for weight management, for unexplained weight loss or poor appetite, for food allergies or intolerances, or for general advice on healthy eating. A doctor can provide a care plan that makes sessions more affordable.

\textbf{Call 000 (or local emergency services) for:}
\begin{itemize}
    \item Severe allergic reactions to food, with difficulty breathing or swelling
    \item Choking on food
    \item Any medical emergency
\end{itemize}

\section*{Diagnosis and Investigations}
\begin{itemize}
    \item \textbf{Dietary assessment:} detailed review of what, when, and how much the person eats and drinks
    \item \textbf{Health assessment:} weight, body measurements, and relevant medical information
    \item \textbf{Nutritional assessment:} review of blood test results and nutritional status
    \item \textbf{Goal setting:} agreement on realistic, achievable targets
    \item \textbf{Follow-up:} monitoring of progress and adjustment of the plan
\end{itemize}

\section*{Management}
\begin{itemize}
    \item \textbf{Individual consultations:} assessment, education, and planning
    \item \textbf{Medical nutrition therapy:} for diabetes, heart disease, kidney disease, and other conditions
    \item \textbf{Weight management:} evidence-based programmes combining diet, activity, and behaviour change
    \item \textbf{Gut health:} advice for irritable bowel syndrome, coeliac disease, and food intolerances
    \item \textbf{Paediatric nutrition:} infant feeding, allergies, and healthy growth
    \item \textbf{Aged care and clinical nutrition:} malnutrition, swallowing difficulties, and tube feeding
    \item \textbf{Referral pathways:} GP chronic disease management plans and private health insurance
\end{itemize}

\section*{Complications}
\begin{itemize}
    \item Poorly managed diet-related conditions without professional guidance
    \item Nutritional deficiencies from fad diets or self-directed restriction
    \item Conflicting and misleading nutrition information online
    \item Cost and access barriers to dietitian services
\end{itemize}

\section*{Prognosis and Outcomes}
Dietitian care produces meaningful improvements in blood glucose, cholesterol, blood pressure, weight, and quality of life for people with diet-related conditions. Most people who attend sessions and make gradual, supported changes achieve lasting results. Individualised advice from a qualified dietitian is more effective and safer than general online guidance.

\section*{Prevention and Self-Care}
\begin{itemize}
    \item See a dietitian early when a condition is diagnosed
    \item Ask your GP about a chronic disease management plan for rebates
    \item Bring your food diary, medications, and blood test results to appointments
    \item Set small, achievable goals between sessions
    \item Be honest about what you eat, so advice can be practical
    \item Check that the practitioner is an accredited practising dietitian
    \item Continue regular follow-up to maintain changes
\end{itemize}

\section*{Sources and Further Reading}
The following reputable sources provide further information for healthcare students and patients seeking reliable, current advice about dietitians and nutrition services.

\begin{itemize}
    \item Dietitians Australia. \textit{Find an accredited practising dietitian.} https://dietitiansaustralia.org.au/
    \item Healthdirect Australia. \textit{Dietitians and nutritionists: what they do.} https://www.healthdirect.gov.au/
    \item Eat For Health (NHMRC). \textit{Australian Dietary Guidelines.} https://www.eatforhealth.gov.au/
    \item Services Australia. \textit{Chronic disease management plans and allied health.} https://www.servicesaustralia.gov.au/
    \item NHS. \textit{Seeing a dietitian.} https://www.nhs.uk/
    \item Dietitian Connection. \textit{Nutrition information and dietitian directory.} https://dietitianconnection.com/
\end{itemize}
